\documentclass[12pt]{article}
\usepackage[utf8]{inputenc}
\usepackage[T1]{fontenc}
\usepackage{hyperref}
\title{Southeastern Pomo Language Resources for K-12 Curriculum Development}
\author{Compiled by Claude Sonnet 5.5 (claude-sonnet-5-5) for Andr\'as Kornai}
\date{}
\begin{document}
\maketitle

\section{Introduction}
Southeastern Pomo (also called Sulfur Bank Pomo, Elem Pomo, Lower Lake Pomo) is a Pomoan language traditionally spoken at the eastern end of Clear Lake, Lake County, California. It is critically endangered. Documentation found in search is modest, and no school programme was found.

\section{Primary Educational Resources}
None found.

\section{Linguistic and Reference Materials}
\begin{itemize}
\item Berkeley Language Center, Southeastern Pomo language page: \url{https://cla.berkeley.edu/languages/southeastern-pomo.html}
\item Southeastern Pomo Language Project (UBC wiki description of the UC Berkeley project and its online dictionary): \url{https://wiki.ubc.ca/Documentation:RelLex/Southeastern_Pomo_Language_Project}
\item Robert Louis Oswalt Papers on Pomoan Languages (archival collection, covering Pomoan languages generally): \url{https://openpolar.no/Record/ftdatacite:10.7297%2Fx2c24tdg}
\end{itemize}
Results on Southern Pomo and Kashaya concern other Pomoan languages and are not included.

\section{Religious and Community Resources}
None found.

\section{Digital and Audio Resources}
The Southeastern Pomo online dictionary is described in the project page above; the dictionary itself was not located as a separate URL in search.

\section{Curriculum Development Context}
A revitalization curriculum would rely on the archival and dictionary materials above and would need to be developed with Lake County Pomo communities. No existing classroom materials were found.
\end{document}